\documentclass[a4paper,11pt]{article}
\usepackage{operationshandbook}
\begin{document}
\begin{center}
{\Large\bfseries 構成・執筆 業務ハンドブック（独立試作）}\par\medskip
{\small 独立設計・試作・内容確認未了}\par
{\small v0.1.0\quad 2026年10月19日}\par
\end{center}
\vspace{6mm}
\noindent 本書は組織から承認を受けた規程又は公式の業務マニュアルを示すものではありません。実務への適用に際しては、必要な法令、利用許諾及び組織の権限規程を確認してください。\par
\tableofcontents
\clearpage
\section{目的、適用範囲と制作工程}
\par 本書は調査・取材資料を読者に届ける制作工程の独立設計案である。原稿、図表、短文、音声・映像台本を対象に、正確性、説明可能性、倫理、情報へのアクセスを同時に検討する。\par\medskip
\par 先行工程で確認した事実、主張、取材発言及び利用条件を継承する。制作過程で新たな疑義が生じれば、調査又は取材へ差し戻す。制作完了は校正、専門確認又は公開承認の成立を意味しない。\par\medskip
\section{文章・図表・媒体に共通する標準}
\par 以下は参照資料を踏まえた独立提案である。学会ガイドライン、政府の公用文作成に関する手引、ウェブ技術規格と本OSの提案にはそれぞれ異なる適用範囲がある。\par\medskip
\subsection{記事構成・読者の理解順序の設計}
\subsubsection{構成設計の目的}
\par 文章の役割、対象読者、主要な問い、中心となる確認済みの事実を先に定める。導入、背景、方法、結果、解釈、限界、今後の課題が果たす説明機能を検討し、媒体の長さと読者に応じて組み立てる。\par\medskip
\subsubsection{見出しと段落}
\par 各段落に中心情報、根拠又は事例、前後とのつながりを持たせる。見出しは本文が答える範囲を超えて断定せず、読者が必要な箇所へ戻れる構造を作る。\par\medskip
\subsubsection{読者への経路}
\par 研究の歴史的経過を含める場合、観測・実験・解釈の時間順と、記事内での説明順を混同しない。関係する事実の発見と報道可能になった時点を分けて記す。\par\medskip
\subsubsection{変更と検証}
\par 編集上の説明順の変更により、因果、時間、対象範囲、引用の趣旨が変わる場合は原資料と照合して修正する。\par\medskip
\subsection{文章・科学的主張の根拠対応}
\subsubsection{原資料との連結}
\par 主張を、出典、原資料の箇所、研究条件、主張の類型（観測、推論、仮説、予測、目標）、原稿内の使用箇所と結び付ける。重要な見出しや図の文言も根拠管理の対象とする。\par\medskip
\subsubsection{表現の強度}
\par 関連があるという研究から因果を断定せず、予備的知見を確立した結論と扱わない。研究対象の人数、母集団、期間、測定条件、効果量・区間推定等の関連情報は重要性に応じて併記する。\par\medskip
\subsubsection{情報の変換}
\par 論文や取材メモを編集用に短縮する際、反対結果、限界、比較対象の省略が意味を変えないか点検する。記事内の説明の単純化によって研究設計の条件を消さない。\par\medskip
\subsubsection{未確認への対応}
\par 根拠が特定できない主張や両立しない資料については、書換えによる解決を試みる前に調査・照会へ戻し、公開に必要な解消条件を記録する。\par\medskip
\subsection{読者に応じた技術日本語・用語運用}
\subsubsection{読者設定}
\par 読む前に求められる知識、専門用語、略語を推定ではなく企画書と対象読者分析に基づき設定する。一般読者向けの説明と専門知識のある読者向けの説明に適した詳細度を選ぶ。\par\medskip
\subsubsection{専門用語の説明}
\par 技術語の初出時は文脈に即した短い説明を置き、必要なら注記やFAQへ接続する。類比はどの特徴を説明し、どこで成立しなくなるか明記する。\par\medskip
\subsubsection{文章上の確認}
\par 一文の対象と述語、修飾関係、主語の変化、時制、固有名詞、助数詞、数値の単位を整える。文章の読みやすさを改善しても原資料の限定条件は維持する。\par\medskip
\subsubsection{用字方針}
\par 媒体別のスタイルブックを参照し、公用文のための政府の一般的な手引を一律の強制規則として扱わない。専門用語と固有名詞の表記方針を用語台帳に保存する。\par\medskip
\subsection{引用・発言・翻訳の編集規律}
\subsubsection{発言の来歴}
\par 直接引用には録音・逐語記録の位置、発言者、日時、前後の文脈を結び付ける。要約と翻訳は直接引用に置き換えず、編集変換の記録を保持する。\par\medskip
\subsubsection{科学的確認}
\par 発言中の技術的な説明は論文等の一次資料で別に確認する。発言が正確に記録されたことと、その内容の科学的妥当性を区分する。\par\medskip
\subsubsection{引用の必要性}
\par 本人の経験や専門的説明を伝える役割を確かめ、過度の切り取りや不均衡な引用が印象を誘導しないか検討する。事前合意した用途と出典表示を確認する。\par\medskip
\subsubsection{疑義の差戻し}
\par 発言に不明確な点があれば取材担当へ照会し、該当箇所を仮置きとして残す。確認済みと未確認の区別が原稿上でも維持されるようにする。\par\medskip
\subsection{図表・キャプション・数量表現の品質}
\subsubsection{出典と再現}
\par 図表ごとにデータ又は原図の出所、利用条件、対象母集団、集計方法、作成者、更新時点を付す。再描画や単位換算は元の数値から追跡可能にする。\par\medskip
\subsubsection{視覚的意味}
\par 目盛、軸の範囲、色と形、分母、ゼロの扱い、差異の強調、統計的不確実性を点検する。視覚上の大小と科学的な重要性を同一とみなさない。\par\medskip
\subsubsection{本文との整合}
\par キャプションに何が示されるか、どの条件下の結果か、出典と限界を簡潔に示し、本文、図、代替説明の主張を相互照合する。\par\medskip
\subsubsection{複数の読み方}
\par 複雑な図には本文で要点を説明し、必要に応じてデータ表、詳細な図解説明、アクセシブルな代替を提供する。\par\medskip
\subsection{媒体横断の情報同等性と設計}
\subsubsection{共通の核}
\par 主要な事実、証拠の確度、重要な制約と関連する引用IDを、文章、SNS、音声、動画、FAQの各媒体に共通する核として設計する。\par\medskip
\subsubsection{媒体固有の構成}
\par 尺や画面の制約を考慮しつつ、研究段階の誤認につながる限定条件の脱落を避ける。動画用台本と文字起こしは対応IDを持つ別バージョンとして管理する。\par\medskip
\subsubsection{任意の利用}
\par 動画や音声を見聞きすることを理解の必須条件にせず、重要情報を適切な代替媒体で参照できるようにする。代替テキスト、字幕、書き起こしの目的を別々に定める。\par\medskip
\subsubsection{整合確認}
\par 各媒体の見出し、数値、引用、公開解禁条件を共通の原資料と照合し、派生版が変更された場合の改訂対象を記録する。\par\medskip
\subsection{AI補助草稿・変換の検証}
\subsubsection{用途}
\par AIは草稿の候補、読みやすい表現案、構成の比較、公開資料の候補整理、用語の問題点抽出等を補助できる。AI出力を科学的根拠、本人発言、著作権の許諾又は公開承認の代替と扱わない。\par\medskip
\subsubsection{入力管理}
\par 公開可能性、権利、個人情報、未公表研究、第三者素材の取扱条件を入力前に確認する。機密資料・録音の外部送信は事前に認められた範囲に限る。\par\medskip
\subsubsection{出力検証}
\par 生成された断定、引用、参考文献、数値、翻訳、人物情報は一次資料や原記録で検証する。根拠を確認できない文を本文に取り込まない。\par\medskip
\subsubsection{過程の記録}
\par AI利用箇所、目的、人による修正、出典との照合、説明責任を引き継げる程度の履歴を残す。\par\medskip
\subsection{制作原稿の版・差戻し・編集引継ぎ}
\subsubsection{文書版}
\par 原稿本文、見出し、図版、字幕、台本、引用表の版を識別し、参照する根拠ID、変更理由、責任者を対応させる。\par\medskip
\subsubsection{差戻しの分類}
\par 新しい科学的な疑問は資料調査へ、発言文脈の不明点は取材へ、出版前の専門的評価は次段階の校閲・確認へ移管する。引継ぎにより問題が自動解決したとみなさない。\par\medskip
\subsubsection{引継ぎの範囲}
\par 完成・未完成の媒体、確認待ちの科学的主張、権利・利用条件、品質上の課題、必要な校正・校閲観点を受け手へ渡す。\par\medskip
\subsubsection{判断の境界}
\par Stage 4の完了は文章の制作上の節目であり、科学的正しさ、法的適合、組織の公開承認を意味しない。\par\medskip
\section{編集倫理、認識上の限界と関係者への配慮}
\par 文章の誇張と表象、読者の多様なアクセス、第三者素材の権利、AI補助の責任を相互に関連づけつつ、それぞれ独立した判断対象として扱う。\par\medskip
\subsection{強調・誇張と誤認を生む編集}
\subsubsection{確度と期待}
\par センセーショナルな見出し、将来の治療・社会応用への過大な期待、検証前の技術を確立した成果として示す文章の影響を評価する。\par\medskip
\subsubsection{省略による偏り}
\par 数字や成功事例を選んで提示する際、対象外の結果、限定条件、失敗や相反する報告の省略が結論を変えるか確認する。\par\medskip
\subsubsection{公共への責任}
\par 読みやすい文章と正確性を同時に追求し、短縮によって誤認が避けられない場合は追加説明への明瞭な経路を提供する。\par\medskip
\subsection{研究者・関係者の尊厳と引用表象}
\subsubsection{当事者の声}
\par 取材対象者の語りや研究参加者の経験を、編集上の目的に合わせて単一の役割に固定しない。文脈、立場、発言の条件を踏まえて掲載する。\par\medskip
\subsubsection{非対称性}
\par 取材側と協力者の権限、知識、公開影響の違いを確認し、修正要請や異議の窓口を明示する。\par\medskip
\subsubsection{少数者と個別事例}
\par 少数者の経験や例外的事例を分野全体の典型と誤認させない。必要に応じて一般化の制約を説明する。\par\medskip
\subsection{情報アクセスと理解の倫理}
\subsubsection{多様な利用者}
\par 専門知識、障害、言語、端末や帯域、時間的制約などによって理解に必要な経路が異なる。多様な媒体で重要情報への選択可能性を確保する。\par\medskip
\subsubsection{アクセシビリティ}
\par 代替テキスト、字幕、見出し構造、用語説明の質を検討し、図版の重要情報が色や音声だけに依存しないようにする。\par\medskip
\subsubsection{合理的な運用}
\par 媒体別の内容品質と負荷を検証し、全媒体が完全に同一となることを目的とせず、重要な意味と条件が保持されることを優先する。\par\medskip
\subsection{制作物・原素材・AIの権利と責任}
\subsubsection{第三者素材}
\par 他者の図表、写真、発言、音楽、書籍の抜粋などは利用許諾、適用される例外、出典、利用範囲を個別に記録する。\par\medskip
\subsubsection{著作者とAI}
\par AIが文章や図版案を作った場合も、使用権、事実確認、人格的利益、責任の所在を人が確認する。出力に含まれる素材の権利が自動確定したとみなさない。\par\medskip
\subsubsection{承認の区別}
\par 専門家による原稿確認、著作物の利用許可、編集責任者による品質判断、公開承認を異なる工程として残す。\par\medskip
\section{取材・調査成果の受領と編集ブリーフの確定}
\subsection{取材・調査成果の受領と編集ブリーフの確定}
\par 担当：編集責任者\par\medskip
\subsubsection*{入力資料}
\begin{enumerate}[leftmargin=2em,itemsep=0.35em]
\item Stage 1企画書
\item Stage 2資料台帳・主張対応表
\item Stage 3取材記録・許諾条件
\end{enumerate}
\subsubsection*{成果物}
\begin{enumerate}[leftmargin=2em,itemsep=0.35em]
\item 編集ブリーフ
\item 利用可能素材一覧
\item 未決リスクの台帳
\end{enumerate}
\subsubsection*{完了の確認事項}
\begin{enumerate}[leftmargin=2em,itemsep=0.35em]
\item 目的、対象読者、参照できる根拠と使用可能な素材が記録された
\item 文書の対象版、編集上の判断、未解決事項が後から追跡できる
\item 担当者と編集責任者が後工程に渡せる範囲と差戻し条件を確認した
\end{enumerate}
\par 人による判断を要する：はい\par\medskip
\subsubsection*{実施手順}
\begin{enumerate}[leftmargin=2em,itemsep=0.65em]
\item 企画目的、対象読者、掲載媒体、原稿の役割と分量を再確認する\par\noindent\textit{記録・成果：}編集ブリーフ\nobreak\smallskip
\item 前工程の資料・取材記録の存在、版、使用条件と担当者を確認する\par\noindent\textit{記録・成果：}利用可能素材一覧\nobreak\smallskip
\item 結果、方法、解釈、予測、研究目標の類型別に使用候補を抽出する\par\noindent\textit{記録・成果：}未決リスクの台帳\nobreak\smallskip
\item 中心的な問いと社会的文脈、研究上の論点、読者の既有知識を記す\par\noindent\textit{記録・成果：}編集ブリーフ\nobreak\smallskip
\item 主要な未確認事項と権利上の未解決事項を制作対象から暫定的に除外する\par\noindent\textit{記録・成果：}利用可能素材一覧\nobreak\smallskip
\item 執筆担当、図表担当、確認窓口、締切と変更経路を整理する\par\noindent\textit{記録・成果：}未決リスクの台帳\nobreak\smallskip
\item 編集責任者が使用可能な根拠、保留素材、構成作業の開始条件を記録する\par\noindent\textit{記録・成果：}編集ブリーフ\nobreak\smallskip
\end{enumerate}
\subsubsection*{例外・差戻し}
\begin{enumerate}[leftmargin=2em,itemsep=0.35em]
\item 調査成果が不足する場合はStage 2に補足を依頼する
\item 原記録・発言の権利条件が不明な場合はStage 3で追加確認する
\end{enumerate}
\section{記事の論理構成・見出しと読者の理解順序}
\subsection{記事の論理構成・見出しと読者の理解順序}
\par 担当：科学編集担当者\par\medskip
\subsubsection*{入力資料}
\begin{enumerate}[leftmargin=2em,itemsep=0.35em]
\item 編集ブリーフ
\item 根拠台帳
\item 研究背景
\item 媒体と対象読者
\end{enumerate}
\subsubsection*{成果物}
\begin{enumerate}[leftmargin=2em,itemsep=0.35em]
\item 構成案
\item 節ごとの中心主張・出典対応表
\item 検証上の留意点
\end{enumerate}
\subsubsection*{完了の確認事項}
\begin{enumerate}[leftmargin=2em,itemsep=0.35em]
\item 目的、対象読者、参照できる根拠と使用可能な素材が記録された
\item 文書の対象版、編集上の判断、未解決事項が後から追跡できる
\item 担当者と編集責任者が後工程に渡せる範囲と差戻し条件を確認した
\end{enumerate}
\par 人による判断を要する：はい\par\medskip
\subsubsection*{実施手順}
\begin{enumerate}[leftmargin=2em,itemsep=0.65em]
\item 問いと確認可能な中心的事実を一文で仮置きする\par\noindent\textit{記録・成果：}構成案\nobreak\smallskip
\item 論文の発表順と読者の理解に必要な説明順の違いを認識する\par\noindent\textit{記録・成果：}節ごとの中心主張・出典対応表\nobreak\smallskip
\item 導入、研究背景、方法、結果、限界、展望の段落機能を定める\par\noindent\textit{記録・成果：}検証上の留意点\nobreak\smallskip
\item 各節に主張IDと根拠ID、未確認の対立見解を付ける\par\noindent\textit{記録・成果：}構成案\nobreak\smallskip
\item 見出しが主張の確度を超えていないか点検する\par\noindent\textit{記録・成果：}節ごとの中心主張・出典対応表\nobreak\smallskip
\item 読者が段落間の関係を理解できる接続文や図の配置を計画する\par\noindent\textit{記録・成果：}検証上の留意点\nobreak\smallskip
\item 編集担当と科学確認担当が構成上の問題点を記録する\par\noindent\textit{記録・成果：}構成案\nobreak\smallskip
\end{enumerate}
\subsubsection*{例外・差戻し}
\begin{enumerate}[leftmargin=2em,itemsep=0.35em]
\item 誤った因果順序を誘導する構成案は修正する
\item 冒頭の印象だけで複雑な結果を単純な成功例に変換しない
\end{enumerate}
\section{研究解説・人物記事の初稿作成}
\subsection{研究解説・人物記事の初稿作成}
\par 担当：科学ライター\par\medskip
\subsubsection*{入力資料}
\begin{enumerate}[leftmargin=2em,itemsep=0.35em]
\item 構成案
\item 一次資料
\item 確認済み引用
\item 読者設定
\end{enumerate}
\subsubsection*{成果物}
\begin{enumerate}[leftmargin=2em,itemsep=0.35em]
\item 初稿
\item 段落別の根拠対応
\item 暫定見出し・脚注
\end{enumerate}
\subsubsection*{完了の確認事項}
\begin{enumerate}[leftmargin=2em,itemsep=0.35em]
\item 目的、対象読者、参照できる根拠と使用可能な素材が記録された
\item 文書の対象版、編集上の判断、未解決事項が後から追跡できる
\item 担当者と編集責任者が後工程に渡せる範囲と差戻し条件を確認した
\end{enumerate}
\par 人による判断を要する：はい\par\medskip
\subsubsection*{実施手順}
\begin{enumerate}[leftmargin=2em,itemsep=0.65em]
\item 各段落の役割と中心となる確認済みの事実を記してから本文を展開する\par\noindent\textit{記録・成果：}初稿\nobreak\smallskip
\item 研究方法、対象、結果、解釈と未解決課題の区分を本文でも保持する\par\noindent\textit{記録・成果：}段落別の根拠対応\nobreak\smallskip
\item 調査で確認した固有名詞、数値、比較条件と出典を追跡できる形で挿入する\par\noindent\textit{記録・成果：}暫定見出し・脚注\nobreak\smallskip
\item 取材発言を逐語・要約・編集上の記述に分けて組み込む\par\noindent\textit{記録・成果：}初稿\nobreak\smallskip
\item 技術語は初出で説明し、類比を用いる場合は成立する範囲を示す\par\noindent\textit{記録・成果：}段落別の根拠対応\nobreak\smallskip
\item 研究の背景・過程・協働とその時間的文脈を必要に応じて加える\par\noindent\textit{記録・成果：}暫定見出し・脚注\nobreak\smallskip
\item 見出し、リード、本文、図の相互整合を通読し不明点を残す\par\noindent\textit{記録・成果：}初稿\nobreak\smallskip
\end{enumerate}
\subsubsection*{例外・差戻し}
\begin{enumerate}[leftmargin=2em,itemsep=0.35em]
\item 事例から集団全体への一般化が難しいときは対象範囲を限定する
\item 原文が確認できない引用は発言として確定しない
\end{enumerate}
\section{研究ニュース・発表原稿・FAQの作成}
\subsection{研究ニュース・発表原稿・FAQの作成}
\par 担当：科学編集担当者\par\medskip
\subsubsection*{入力資料}
\begin{enumerate}[leftmargin=2em,itemsep=0.35em]
\item 一次資料
\item 研究解説草稿
\item 発表状況
\item 媒体別の分量方針
\end{enumerate}
\subsubsection*{成果物}
\begin{enumerate}[leftmargin=2em,itemsep=0.35em]
\item 短いニュース草稿
\item 発表原稿案
\item FAQと保留事項
\end{enumerate}
\subsubsection*{完了の確認事項}
\begin{enumerate}[leftmargin=2em,itemsep=0.35em]
\item 目的、対象読者、参照できる根拠と使用可能な素材が記録された
\item 文書の対象版、編集上の判断、未解決事項が後から追跡できる
\item 担当者と編集責任者が後工程に渡せる範囲と差戻し条件を確認した
\end{enumerate}
\par 人による判断を要する：はい\par\medskip
\subsubsection*{実施手順}
\begin{enumerate}[leftmargin=2em,itemsep=0.65em]
\item 短文の目的、対象読者、公開状況、記事の文脈を確定する\par\noindent\textit{記録・成果：}短いニュース草稿\nobreak\smallskip
\item 何が確かめられたか、対象範囲、研究の段階と重要な限定条件を整理する\par\noindent\textit{記録・成果：}発表原稿案\nobreak\smallskip
\item タイトル・リードに観測と推測の境界を保持する\par\noindent\textit{記録・成果：}FAQと保留事項\nobreak\smallskip
\item プレス向け原稿の発表主体、論文出典、研究手法、共同研究情報と問合せ先を確認する\par\noindent\textit{記録・成果：}短いニュース草稿\nobreak\smallskip
\item FAQは読者の疑問を独立した問いとして構成し、過度に断定せず回答する\par\noindent\textit{記録・成果：}発表原稿案\nobreak\smallskip
\item 用途別の分量制限による意味の脱落を比較表に記録する\par\noindent\textit{記録・成果：}FAQと保留事項\nobreak\smallskip
\item 内容と公開条件を編集責任者が確認する\par\noindent\textit{記録・成果：}短いニュース草稿\nobreak\smallskip
\end{enumerate}
\subsubsection*{例外・差戻し}
\begin{enumerate}[leftmargin=2em,itemsep=0.35em]
\item 未発表結果の発信は公開解禁条件と研究者の確認へ戻す
\item 短文に必要な留保が入らない場合は詳細説明先の明示を検討する
\end{enumerate}
\section{専門表現・日本語の可読性と用語統一}
\subsection{専門表現・日本語の可読性と用語統一}
\par 担当：日本語編集担当者\par\medskip
\subsubsection*{入力資料}
\begin{enumerate}[leftmargin=2em,itemsep=0.35em]
\item 各種草稿
\item 用語台帳
\item 対象読者設定
\item 原資料
\end{enumerate}
\subsubsection*{成果物}
\begin{enumerate}[leftmargin=2em,itemsep=0.35em]
\item 表現改訂案
\item 用語・略語台帳
\item 難解事項一覧
\end{enumerate}
\subsubsection*{完了の確認事項}
\begin{enumerate}[leftmargin=2em,itemsep=0.35em]
\item 目的、対象読者、参照できる根拠と使用可能な素材が記録された
\item 文書の対象版、編集上の判断、未解決事項が後から追跡できる
\item 担当者と編集責任者が後工程に渡せる範囲と差戻し条件を確認した
\end{enumerate}
\par 人による判断を要する：はい\par\medskip
\subsubsection*{実施手順}
\begin{enumerate}[leftmargin=2em,itemsep=0.65em]
\item 文体、敬体・常体、表記の方針を媒体別に決める\par\noindent\textit{記録・成果：}表現改訂案\nobreak\smallskip
\item 専門用語と略語を初出と既出で区分し読者に必要な説明を整える\par\noindent\textit{記録・成果：}用語・略語台帳\nobreak\smallskip
\item 主語、述語、修飾先、論理関係、時制と代名詞の指す対象を確認する\par\noindent\textit{記録・成果：}難解事項一覧\nobreak\smallskip
\item 概念の言い換えによって比較条件や科学的限定が脱落しないか原資料と照合する\par\noindent\textit{記録・成果：}表現改訂案\nobreak\smallskip
\item 数値、単位、地名、人名、所属、用字を台帳と整合させる\par\noindent\textit{記録・成果：}用語・略語台帳\nobreak\smallskip
\item 複数の読者層向けに説明の深さを変更する場合は科学的な核を維持する\par\noindent\textit{記録・成果：}難解事項一覧\nobreak\smallskip
\item 改訂履歴と科学的確認が必要な言い換えを明示する\par\noindent\textit{記録・成果：}表現改訂案\nobreak\smallskip
\end{enumerate}
\subsubsection*{例外・差戻し}
\begin{enumerate}[leftmargin=2em,itemsep=0.35em]
\item わかりやすい表現が科学的事実を変える場合は事実確認へ戻す
\item 新たな用語の訳語が定着していない場合は暫定表記の根拠を付す
\end{enumerate}
\section{引用・科学的主張・叙述の対応確認}
\subsection{引用・科学的主張・叙述の対応確認}
\par 担当：科学編集担当者\par\medskip
\subsubsection*{入力資料}
\begin{enumerate}[leftmargin=2em,itemsep=0.35em]
\item 取材逐語記録
\item 引用台帳
\item 主張台帳
\item 各種初稿
\end{enumerate}
\subsubsection*{成果物}
\begin{enumerate}[leftmargin=2em,itemsep=0.35em]
\item 文章と原資料の対応表
\item 要照会文言と修正記録
\end{enumerate}
\subsubsection*{完了の確認事項}
\begin{enumerate}[leftmargin=2em,itemsep=0.35em]
\item 目的、対象読者、参照できる根拠と使用可能な素材が記録された
\item 文書の対象版、編集上の判断、未解決事項が後から追跡できる
\item 担当者と編集責任者が後工程に渡せる範囲と差戻し条件を確認した
\end{enumerate}
\par 人による判断を要する：はい\par\medskip
\subsubsection*{実施手順}
\begin{enumerate}[leftmargin=2em,itemsep=0.65em]
\item 本文中の引用文、翻訳、意訳、記者による推論を別々に識別する\par\noindent\textit{記録・成果：}文章と原資料の対応表\nobreak\smallskip
\item 直接引用は時点・発言者・原音・逐語の対応を照合する\par\noindent\textit{記録・成果：}要照会文言と修正記録\nobreak\smallskip
\item 数字と研究主張を一次資料の方法・対象・結論に対応づける\par\noindent\textit{記録・成果：}文章と原資料の対応表\nobreak\smallskip
\item 見出し・脚注・図版でも因果の断定や条件の欠落がないか調べる\par\noindent\textit{記録・成果：}要照会文言と修正記録\nobreak\smallskip
\item 取材時の利用条件、第三者著作物、未公開情報と公開方法を確認する\par\noindent\textit{記録・成果：}文章と原資料の対応表\nobreak\smallskip
\item 論争的な箇所や発言の意味が変わる箇所を追加照会として登録する\par\noindent\textit{記録・成果：}要照会文言と修正記録\nobreak\smallskip
\item 引用管理台帳の状態と変更履歴を更新する\par\noindent\textit{記録・成果：}文章と原資料の対応表\nobreak\smallskip
\end{enumerate}
\subsubsection*{例外・差戻し}
\begin{enumerate}[leftmargin=2em,itemsep=0.35em]
\item 研究者が記事の事実確認をしたことを本文全体の科学的保証と扱わない
\item 引用の条件に疑義があれば掲載判断を保留する
\end{enumerate}
\section{図表・写真・キャプションと説明資料の設計}
\subsection{図表・写真・キャプションと説明資料の設計}
\par 担当：図表制作担当者と科学編集担当者\par\medskip
\subsubsection*{入力資料}
\begin{enumerate}[leftmargin=2em,itemsep=0.35em]
\item データ・原図
\item 図表の目的
\item 利用条件
\item 草稿本文
\end{enumerate}
\subsubsection*{成果物}
\begin{enumerate}[leftmargin=2em,itemsep=0.35em]
\item 図表仕様書
\item キャプション
\item 出典・権利記録
\item 代替説明
\end{enumerate}
\subsubsection*{完了の確認事項}
\begin{enumerate}[leftmargin=2em,itemsep=0.35em]
\item 目的、対象読者、参照できる根拠と使用可能な素材が記録された
\item 文書の対象版、編集上の判断、未解決事項が後から追跡できる
\item 担当者と編集責任者が後工程に渡せる範囲と差戻し条件を確認した
\end{enumerate}
\par 人による判断を要する：はい\par\medskip
\subsubsection*{実施手順}
\begin{enumerate}[leftmargin=2em,itemsep=0.65em]
\item 各図の目的と読者が理解すべき要点を一つずつ定める\par\noindent\textit{記録・成果：}図表仕様書\nobreak\smallskip
\item データの出典、定義、対象、算出方法、時点と利用権限を記録する\par\noindent\textit{記録・成果：}キャプション\nobreak\smallskip
\item 目盛、軸、単位、分母、縮尺、表示順序、誤差の表現を設計する\par\noindent\textit{記録・成果：}出典・権利記録\nobreak\smallskip
\item グラフや図による関係の可視化と因果の実証の違いを確認する\par\noindent\textit{記録・成果：}代替説明\nobreak\smallskip
\item 本文との整合、キャプションの限定条件と出典の表示を点検する\par\noindent\textit{記録・成果：}図表仕様書\nobreak\smallskip
\item 情報を運ぶ図の代替テキスト・長い説明又はデータ表を準備する\par\noindent\textit{記録・成果：}キャプション\nobreak\smallskip
\item 改訂の経緯、最終確認者、再利用範囲を記録する\par\noindent\textit{記録・成果：}出典・権利記録\nobreak\smallskip
\end{enumerate}
\subsubsection*{例外・差戻し}
\begin{enumerate}[leftmargin=2em,itemsep=0.35em]
\item 画像の見た目で効果が誇張される場合は軸・尺度を改める
\item 図の権利が不明なら公開向け制作を保留する
\end{enumerate}
\section{音声・映像・SNS・アクセシブルな説明の展開}
\subsection{音声・映像・SNS・アクセシブルな説明の展開}
\par 担当：マルチモーダル制作担当者\par\medskip
\subsubsection*{入力資料}
\begin{enumerate}[leftmargin=2em,itemsep=0.35em]
\item 本文草稿
\item 図表仕様
\item 公開可能素材
\item 対象媒体の要件
\end{enumerate}
\subsubsection*{成果物}
\begin{enumerate}[leftmargin=2em,itemsep=0.35em]
\item 動画台本
\item 音声台本
\item 短文原稿
\item 字幕・代替テキスト
\end{enumerate}
\subsubsection*{完了の確認事項}
\begin{enumerate}[leftmargin=2em,itemsep=0.35em]
\item 目的、対象読者、参照できる根拠と使用可能な素材が記録された
\item 文書の対象版、編集上の判断、未解決事項が後から追跡できる
\item 担当者と編集責任者が後工程に渡せる範囲と差戻し条件を確認した
\end{enumerate}
\par 人による判断を要する：はい\par\medskip
\subsubsection*{実施手順}
\begin{enumerate}[leftmargin=2em,itemsep=0.65em]
\item 媒体の目的と利用者の選択可能性を確定する\par\noindent\textit{記録・成果：}動画台本\nobreak\smallskip
\item 主要な科学的事実、制約、引用、出典を全媒体で共有する核として定義する\par\noindent\textit{記録・成果：}音声台本\nobreak\smallskip
\item 動画の場面、音声の語り、SNSの短文を個別に設計する\par\noindent\textit{記録・成果：}短文原稿\nobreak\smallskip
\item 視覚や音声だけに依存する重要事項にはテキスト等の代替を準備する\par\noindent\textit{記録・成果：}字幕・代替テキスト\nobreak\smallskip
\item 時間長や文字数によって研究上の留保が落ちていないか媒体間で比較する\par\noindent\textit{記録・成果：}動画台本\nobreak\smallskip
\item 字幕、書き起こし、画面説明、図版の権利と個人情報を点検する\par\noindent\textit{記録・成果：}音声台本\nobreak\smallskip
\item 担当者が内容の差と整合性を記録して専門確認へ渡す\par\noindent\textit{記録・成果：}短文原稿\nobreak\smallskip
\end{enumerate}
\subsubsection*{例外・差戻し}
\begin{enumerate}[leftmargin=2em,itemsep=0.35em]
\item 動画未視聴でも重要情報を得られる経路を用意する
\item 音声・字幕に原稿と異なる科学的断定があれば改訂する
\end{enumerate}
\section{制作原稿の整合確認・版管理と次段階への引継ぎ}
\subsection{制作原稿の整合確認・版管理と次段階への引継ぎ}
\par 担当：編集責任者\par\medskip
\subsubsection*{入力資料}
\begin{enumerate}[leftmargin=2em,itemsep=0.35em]
\item 本文・見出し・図表・台本
\item 根拠台帳
\item 引用台帳
\item 変更記録
\end{enumerate}
\subsubsection*{成果物}
\begin{enumerate}[leftmargin=2em,itemsep=0.35em]
\item 制作物の版セット
\item 未解決事項台帳
\item 編集・校閲への引継ぎ
\end{enumerate}
\subsubsection*{完了の確認事項}
\begin{enumerate}[leftmargin=2em,itemsep=0.35em]
\item 目的、対象読者、参照できる根拠と使用可能な素材が記録された
\item 文書の対象版、編集上の判断、未解決事項が後から追跡できる
\item 担当者と編集責任者が後工程に渡せる範囲と差戻し条件を確認した
\end{enumerate}
\par 人による判断を要する：はい\par\medskip
\subsubsection*{実施手順}
\begin{enumerate}[leftmargin=2em,itemsep=0.65em]
\item 必要な媒体と文書の有無、版番号、担当、状態を一覧化する\par\noindent\textit{記録・成果：}制作物の版セット\nobreak\smallskip
\item 主張・出典ID、図表と本文、引用と原記録の対応を確認する\par\noindent\textit{記録・成果：}未解決事項台帳\nobreak\smallskip
\item 誇張、重大な権利疑義、誤訳、欠けた限定条件と未回答事項を抽出する\par\noindent\textit{記録・成果：}編集・校閲への引継ぎ\nobreak\smallskip
\item 未確認事項の重要性に応じ差戻し先、担当と期限を決める\par\noindent\textit{記録・成果：}制作物の版セット\nobreak\smallskip
\item 後続の編集・校正・校閲が点検すべき重点箇所を示す\par\noindent\textit{記録・成果：}未解決事項台帳\nobreak\smallskip
\item 根拠が確認された箇所と本文上未完成の箇所を分けて引き渡す\par\noindent\textit{記録・成果：}編集・校閲への引継ぎ\nobreak\smallskip
\item 編集責任者が制作上の節目を確認し、公開承認の未実施を明示する\par\noindent\textit{記録・成果：}制作物の版セット\nobreak\smallskip
\end{enumerate}
\subsubsection*{例外・差戻し}
\begin{enumerate}[leftmargin=2em,itemsep=0.35em]
\item 重大な根拠不足はStage 2へ差し戻す
\item 発言の利用条件の疑義はStage 3へ戻す
\item 本段階の完了を公開承認とみなさない
\end{enumerate}
\section{工程内のチェックリスト}
\par チェックリストへの記入だけで研究の妥当性、利用権限又は公開承認は成立しない。実案件では状態、根拠、担当、日時、異議、改訂をChecklistRunとして保持する。\par\medskip
\subsection{科学的不確実性・誇張防止チェックリスト}
\par 確認用の定義であり、個別案件での実施状況、根拠資料、確認者、レビューの結果は別途保存する。\par\medskip
\subsubsection*{確認項目}
\begin{enumerate}[leftmargin=2em,itemsep=0.3em]
\item 本文・見出し・図における「観測、推論、仮説、予測、目標」を適切に区分した\par
\item 相関から因果を断定していない\par
\item 研究対象、人数、期間、比較対象と適用範囲を必要な箇所に示した\par
\item 数値の見せ方に効果量や不確実性の重要な欠落がない\par
\item 研究結果から臨床・社会実装までの距離を過度に縮めていない\par
\item 発表の新規性と既存研究の連続性・異なる解釈を検討した\par
\item タイトル・短文・図表の表現が原稿本文より過度に強くなっていない\par
\item 検証が必要な表現には主張ID、調査担当と修正期限を付した\par
\end{enumerate}
\subsection{アクセシビリティ品質チェックリスト}
\par 確認用の定義であり、個別案件での実施状況、根拠資料、確認者、レビューの結果は別途保存する。\par\medskip
\subsubsection*{確認項目}
\begin{enumerate}[leftmargin=2em,itemsep=0.3em]
\item 情報を運ぶ画像に目的に応じた代替テキストが存在する\par
\item 複雑なグラフに主要な数値・傾向を説明する代替経路がある\par
\item 音声・動画に必要な字幕・書き起こしと説明文を用意した\par
\item 見出し構造と文書の読み順が内容の関係を表している\par
\item 色、音、動きの一つの感覚手段だけに重要情報を依存させていない\par
\item 専門語・略語の説明が対象読者に適切な位置にある\par
\item 異なる端末・表示環境で読解不能な部分がないか確認する\par
\item 確認者、参照した技術基準、未解決の課題を記録した\par
\end{enumerate}
\subsection{複数媒体間の内容整合チェックリスト}
\par 確認用の定義であり、個別案件での実施状況、根拠資料、確認者、レビューの結果は別途保存する。\par\medskip
\subsubsection*{確認項目}
\begin{enumerate}[leftmargin=2em,itemsep=0.3em]
\item 長文、ニュース、FAQ、SNS、動画、音声で中心的な主張が一致する\par
\item 対象となる人数、期間、適用条件の記載に重要な矛盾がない\par
\item 直接引用、要約、翻訳に同じ原記録IDが付いている\par
\item 画像・キャプション・字幕が本文より強く断定していない\par
\item 短縮版で除かれた限定条件へ読者が参照できる経路がある\par
\item 媒体ごとの出典と権利、解禁条件を確認した\par
\item ある媒体の改訂が他媒体へ反映されたか確認した\par
\item 版・確認担当・未決項目を媒体別に残した\par
\end{enumerate}
\subsection{記事構成・論理確認チェックリスト}
\par 確認用の定義であり、個別案件での実施状況、根拠資料、確認者、レビューの結果は別途保存する。\par\medskip
\subsubsection*{確認項目}
\begin{enumerate}[leftmargin=2em,itemsep=0.3em]
\item 記事の目的、読者と中心となる問いが明確である\par
\item 見出し、導入、方法、結果、解釈、限界の各段落に説明上の役割がある\par
\item 重要な主張・数値・引用の箇所に出典IDを結び付けた\par
\item 文章の構成が時間順と因果の区別を壊していない\par
\item 異なる解釈や研究上の限界を必要な程度扱った\par
\item 各節の結論が根拠の確実性と適用範囲を超えない\par
\item 冒頭、見出し、結語と本文の確度が一致する\par
\item 編集上の未決事項と戻り先を記録した\par
\end{enumerate}
\subsection{わかりやすい日本語チェックリスト}
\par 確認用の定義であり、個別案件での実施状況、根拠資料、確認者、レビューの結果は別途保存する。\par\medskip
\subsubsection*{確認項目}
\begin{enumerate}[leftmargin=2em,itemsep=0.3em]
\item 対象読者に必要な技術語を初出で説明した\par
\item 略語、固有名詞、研究段階と専門概念を一貫して表記した\par
\item 一文の主語・述語・修飾先、指示語の対象が理解できる\par
\item 数字、時点、単位、比較対象が読み誤られにくい\par
\item 類比を使った箇所で成立する範囲と限界を説明した\par
\item 読みやすくする変更で因果・適用範囲・留保が変化していない\par
\item 媒体に適切な句読点、見出し、段落の長さを確認した\par
\item 科学的意味が変わり得る言い換えを専門確認へ渡した\par
\end{enumerate}
\subsection{図表・キャプション確認チェックリスト}
\par 確認用の定義であり、個別案件での実施状況、根拠資料、確認者、レビューの結果は別途保存する。\par\medskip
\subsubsection*{確認項目}
\begin{enumerate}[leftmargin=2em,itemsep=0.3em]
\item 図表の原データ又は原図、版、出典、権利状況を記録した\par
\item 軸、単位、分母、サンプル数、比較条件が正確である\par
\item 縮尺、目盛範囲、色と形、注記に誘導的な表現がない\par
\item 有意差・効果量・不確実性の重要な情報が図と本文で一致する\par
\item キャプションが図の対象、結論の射程、出典を説明する\par
\item 本文の記述と図が食い違っていない\par
\item 画像の内容に応じた代替テキスト又は詳細説明を準備した\par
\item 最終データ更新日とレビュー担当を記録した\par
\end{enumerate}
\subsection{映像・音声台本確認チェックリスト}
\par 確認用の定義であり、個別案件での実施状況、根拠資料、確認者、レビューの結果は別途保存する。\par\medskip
\subsubsection*{確認項目}
\begin{enumerate}[leftmargin=2em,itemsep=0.3em]
\item 動画・音声の構成が原稿本文と出典に対応している\par
\item ナレーション、画面文字、テロップ・字幕の数値と固有名詞が一致する\par
\item 見せ方や効果音によって研究結果を過度に強調していない\par
\item 引用音声、人物、画像、音楽等の利用条件を確認した\par
\item 字幕・書き起こし・図の代替説明を準備した\par
\item 再生しなくても重要な情報へアクセスできる経路がある\par
\item 媒体ごとの解禁条件、リンク先、版を揃えた\par
\item 外部制作物の修正責任と未決事項を記録した\par
\end{enumerate}
\section{制作の引継ぎ、品質確認と差戻し}
\par 本文・図表・台本の対象版、主張・根拠ID、引用と原記録の対応、権利・解禁条件、未確認の科学的主張、校閲上の重点項目を次工程に引き渡す。\par\medskip
\par 制作上の確認は次の編集・校正・校閲、専門確認、公開承認から独立して記録する。科学的疑義は調査、取材発言の疑義は取材へ戻し、停止の判断と理由を残す。\par\medskip
\section{参照資料・制度上の射程と改訂}
\par 参照資料の正式な所在と適用上の留意事項は content/operational/stage\_04/REFERENCES.md に記す。研究成果のプレスリリース指針は一般的な独立報道の全工程を規定せず、政府の公用文作成の手引も全媒体を法的に拘束するものではない。\par\medskip
\par 規格とプロセス案は版・確認日・文書IDで参照し、改訂を追跡する。電子的に閲覧する文書と印刷物では代替テキスト等の機能面が異なるため、媒体固有の品質確認を別に行う。\par\medskip
\end{document}
